\documentclass[10pt,a4paper]{article}
\usepackage{amsmath,amssymb,amsthm,graphicx}
\usepackage{hyperref}
\usepackage{fontspec}
\title{Transfer of Reed-Muller local testing to multiplicity codes with quantitative soundness}
\author{Dan Karliner, Roie Salama, Amnon Ta-Shma}
\date{doi:10.4230/LIPIcs.CCC.2022.14; CC BY 4.0}
\begin{document}
\maketitle
\noindent\textit{Editorial scope.} Counted original Theorem51 (introduced as Theorem4): quantitative transfer from an assumed RM k-dimensional tester to the MRM tester, with its exact alpha-prime coefficient and parameter restrictions. Full Claims53–56 and final derivative reconstruction are included. No support claim or numerical corollary is counted separately. Original author mathematics is retained as published.\par\medskip
\noindent\textit{Allowed standard prerequisites (editorial).}\begin{itemize}
\item Reed-Muller distance bound [10], standard Hasse derivative restriction identity [15], Markov and Schwartz-Zippel bounds. Author Lemmas11,15,17 and exact alphabet-reduction definition48 retained. Main statement retains the assumed RM tester soundness instead of claiming it unconditionally.
\end{itemize}\medskip
% Faithful official CCC2022 article14 transcription: necessary definitions and complete Section5 transfer theorem.
\section*{2 Preliminaries (author context excerpts)}
\subsection*{2.1 Notation}
We denote vectors by bold letters. For $\boldsymbol X=X_1,\ldots,X_m$ we denote by $\mathbb F[\boldsymbol X]$ the set of multivariate polynomials in the variables $X_1,\ldots,X_m$. We denote by $\mathbb F[\boldsymbol X]^{\leq d}$ the set of polynomials of individual degree at most $d$, and by $\mathbb F[\boldsymbol X]^{\mathrm{loc}\leq d}$ the set of polynomials of individual degree at most $d$ (i.e degree in each variable). Given a vector $\boldsymbol I=(i_1,\ldots,i_m)\in\mathbb N^m$, we use the notation
\[\boldsymbol X^{\boldsymbol I}\stackrel{\mathrm{def}}=\prod_{k=1}^mX_k^{i_k}.\]
\subsection*{2.2 Reed Muller and Multiplicity codes}
\paragraph{Definition 10.} Let $d,m$ be non-negative integers, and $q$ a prime power. The $(m,d,q)$ Reed-Muller code is defined as the set of evaluation vectors of $m$-variate polynomials of degree $\leq d$ over $\mathbb F_q^m$, namely,
\[\mathrm{RM}(q,m,d)=\{(f(\alpha))_{\alpha\in\mathbb F_q^m}\mid f\in\mathbb F_q[\boldsymbol X]^{\leq d}\}.\]
We will make use of the following lemma:
\paragraph{Lemma 11 ([10], Lemma 3.2).} Let $\delta_{q,m,d}^{\mathrm{RM}}$ be the relative distance of the code $\mathrm{RM}(q,m,d)$. Then $\delta_{q,m,d}^{\mathrm{RM}}\geq q^{-d/(q-1)}$.
\paragraph{Definition 12 (Hasse derivative).} For a multivariate $P(\boldsymbol X)\in\mathbb F[\boldsymbol X]$ where $\boldsymbol X=(X_1,\ldots,X_m)$ for some $m\in\mathbb N$, and a non-negative vector $\boldsymbol I\in\mathbb N^m$, the $\boldsymbol I$-th Hasse derivative of $P$, denoted $P^{(\boldsymbol I)}(\boldsymbol X)$, is the coefficient of $\boldsymbol Z^{\boldsymbol I}$ in the polynomial $P(\boldsymbol X,\boldsymbol Z)=P(\boldsymbol X+\boldsymbol Z)$. Thus
\[P(\boldsymbol X+\boldsymbol Z)=\sum_{\boldsymbol I}P^{(\boldsymbol I)}(\boldsymbol X)\cdot\boldsymbol Z^{\boldsymbol i}.\]
Hasse derivatives are linear. I.e, for all $P,Q\in\mathbb F[\boldsymbol X]$ and $\lambda\in\mathbb F$, $(\lambda P)^{(\boldsymbol I)}(\boldsymbol X)=\lambda P(\boldsymbol X)$ and $P^{(\boldsymbol I)}(\boldsymbol X)+Q^{(\boldsymbol I)}(\boldsymbol X)=(P+Q)^{(\boldsymbol I)}(\boldsymbol X)$. The product rule shows $(PQ)^{(\boldsymbol I)}(\boldsymbol X)=\sum_{\boldsymbol I_0+\boldsymbol I_1=\boldsymbol I}P^{(\boldsymbol I_0)}(\boldsymbol X)\cdot Q^{\boldsymbol I_1}(\boldsymbol X)$.
\paragraph{Definition 13 (Weight).} If $\boldsymbol I=(i_1,\ldots,i_m)\in\mathbb N^m$ then $w(\boldsymbol I)=\sum_{j=1}^mi_j$.
\paragraph{Definition 14 (Multiplicity).} For $P(\boldsymbol X)\in\mathbb F[\boldsymbol X]$ and $\boldsymbol a\in\mathbb F^m$, the multiplicity of $P$ at $\boldsymbol a$, denoted $\operatorname{mult}(P,\boldsymbol a)$, is the largest integer $s$ such that for every non-negative vector $\boldsymbol I$ with $w(\boldsymbol I)<s$ we have $P^{(\boldsymbol I)}(\boldsymbol a)=0$. If $s$ may be taken arbitrarily large, we set $\operatorname{mult}(P,\boldsymbol a)=\infty$.
Note that by definition $\operatorname{mult}(P,\boldsymbol a)\geq0$ for every $\boldsymbol a$. One important property about multiplicities is a generalization of the Schwartz-Zippel lemma for multivariate polynomials:
\paragraph{Lemma 15 ([6]).} Let $P\in\mathbb F[\boldsymbol X]$ be a non-zero polynomial of total degree at most $d$. Then for any finite $A\subseteq\mathbb F$,
\[\sum_{\boldsymbol a\in A^m}\operatorname{mult}(P,\boldsymbol a)\leq d\cdot|A|^{m-1}.\]
\paragraph{Definition 16 (Multiplicity code).} Let $m,d\geq s$ be non-negative integers, and let $q$ be a prime power. Let
\[\Sigma_{m,s}=\mathbb F_q^{\{\boldsymbol I:w(\boldsymbol I)<s\}}\simeq\mathbb F_q^{\binom{m+s-1}{m}}.\]
For $P(\boldsymbol X)\in\mathbb F_q[X_1,\ldots,X_m]$ we define the order $s$ evaluation of $P$ at $\boldsymbol a$, denoted $P^{(<s)}(\boldsymbol a)$, to be the vector $(P^{(\boldsymbol I)}(\boldsymbol a))_{\boldsymbol I:w(\boldsymbol I)<s}\in\Sigma_{m,s}$. The multiplicity code $\mathrm{MRM}(q,m,d,s)$ is defined as follows. The alphabet of the code is $\Sigma_{m,s}$ and the length is $q^m$. Every polynomial $P(\boldsymbol X)\in\mathbb F_q[\boldsymbol X]$ of $\deg(P)\leq d$ defines a codeword by $(P^{(<s)}(\boldsymbol a))_{\boldsymbol a:\boldsymbol a\in\mathbb F_q^m}\in(\Sigma_{m,s})^{q^m}$.
We also let $\mathrm{MRS}(q,d,s):=\mathrm{Mult}(q,1,d,s)$ stand for Reed-Solomon multiplicity code.
The following lemma states the relationship between the derivatives of a polynomial to the derivatives of its restriction to a line.
\paragraph{Lemma 17 ([15], Sec 4).} Let $P\in\mathbb F[\boldsymbol X]$ be a multivariate polynomial where $\boldsymbol X=(X_1,\ldots,X_m)$. Let $\boldsymbol a,\boldsymbol b\in\mathbb F^m$ and define a univariate polynomial $P|_{L_{\boldsymbol a,\boldsymbol b}}(t)=P(\boldsymbol a t+\boldsymbol b)$. Then
\[(P|_{L_{\boldsymbol a,\boldsymbol b}})^{(j)}(t)=\sum_{\boldsymbol I:w(\boldsymbol I)=j}P^{(\boldsymbol I)}(\boldsymbol a t+\boldsymbol b)\cdot\boldsymbol a^{\boldsymbol I}.\]
\section*{3 Polynomials and tables (author context excerpts)}
A table is an element of $\Sigma_{m,s}^{q^m}$, i.e., a function mapping every evaluation point in $\mathbb F_q^m$ to an element in $\Sigma_{m,s}$. EVAL takes a multi-variate polynomial and returns its table of evaluations. More precisely,
\paragraph{Definition 39.} We define $\mathrm{EVAL}_{m,s}:\mathbb F_q[\boldsymbol X]\times\mathbb F_q^m\to\Sigma_{m,s}$ by
\[\mathrm{EVAL}_{m,s}(P;\boldsymbol a)=(P^{(\boldsymbol I)}(\boldsymbol a))_{w(\boldsymbol I)<s}.\]
Similarly, we define $\mathrm{EVAL}_{m,s}:\mathbb F_q[\boldsymbol X]\to(\Sigma_{m,s})^{q^m}$ by
\[\mathrm{EVAL}_{m,s}(P)=(\mathrm{EVAL}(P;\boldsymbol a))_{\boldsymbol a\in\mathbb F_q^m}.\]
We say $T\in(\Sigma_{m,s})^{q^m}$ is the table of $P\in\mathbb F_q[\boldsymbol X]$ if $\mathrm{EVAL}(P)=T$.
An element in $\Sigma_{m,s}$ contains information about all the derivatives of order up to $s$. Sometimes we would like to focus on a certain directional derivative. If $\sigma\in\Sigma_{m,s}$ and $\boldsymbol I\in\mathbb N^m$ with $w(\boldsymbol I)<s$, then $\sigma_{(\boldsymbol I)}\in\mathbb F_q$ is the entry of $\sigma$ that encodes the $\boldsymbol I$'th derivative. Similarly, if $f\in\Sigma_{m,s}^{q^m}$ is a table, i.e., $f:\mathbb F_q^m\to\Sigma_{m,s}$, then $f_{(\boldsymbol I)}:\mathbb F_q^m\to\mathbb F_q$ is defined by letting $f_{(\boldsymbol I)}(\boldsymbol x)$ be the $\boldsymbol I$'th entry of $f(\boldsymbol x)\in\Sigma_{m,s}$.
To understand the low-dimensional tests on tables, we need to define the restriction of tables to subspaces. If $T$ is a table $T:\mathbb F_q^m\to\Sigma_{m,s}$ and we want to restrict to the plane $\boldsymbol a t+\boldsymbol b r+\boldsymbol c$ the restriction $T_{\boldsymbol a,\boldsymbol b,\boldsymbol c}$ should be a table $\mathbb F_q^2\to\Sigma_{2,s}$. To this end we define the alphabet reduction map $\phi_{(\boldsymbol a,\boldsymbol b)}:\Sigma_{m,s}\to\Sigma_{2,s}$ by
\paragraph{Definition 48.}
\[(\phi_{(\boldsymbol a,\boldsymbol b)})(\boldsymbol z)_{\boldsymbol J=(j_1,j_2)}=\sum_{\boldsymbol I\in\mathbb N^m}z_{\boldsymbol I}\cdot\sum_{\substack{\boldsymbol I_1+\boldsymbol I_2=\boldsymbol I\\w(\boldsymbol I_1)=j_1,w(\boldsymbol I_2)=j_2}}\binom{\boldsymbol I}{\boldsymbol I_1}\boldsymbol a^{\boldsymbol I_1}\boldsymbol b^{\boldsymbol I_2}.\]
This map applies the ``chain rule'' to an element in $\Sigma_{m,s}$, in accordance with Lemma 18. We may then define
\[T_{\boldsymbol a,\boldsymbol b,\boldsymbol c}=\phi_{(\boldsymbol a,\boldsymbol b)}\circ T\circ\ell_{\boldsymbol a,\boldsymbol b,\boldsymbol c}.\]
Similarly, for $\boldsymbol h=(\boldsymbol h_0,\boldsymbol h_1,\ldots,\boldsymbol h_k)$ we define $\phi_{\boldsymbol h}$ and $T_{\boldsymbol h}$ by:
\[(\phi_{\boldsymbol h}(\boldsymbol z))_{\boldsymbol J}=\sum_{\boldsymbol I\in\mathbb N^m}z_{\boldsymbol I}\cdot\sum_{\substack{\boldsymbol I_1+\cdots+\boldsymbol I_k=\boldsymbol I\\w(\boldsymbol I_r)=j_r}}\binom{\boldsymbol I}{\boldsymbol I_1,\ldots,\boldsymbol I_k}\prod_{i=1}^k\boldsymbol h_k^{\boldsymbol I_k},\text{ and,}\]
\[T_{\boldsymbol h}(\boldsymbol Y)=\phi_{(\boldsymbol h)}\circ T\left(\boldsymbol h_0+\sum_{i=1}^kY_i\boldsymbol h_i\right).\]
\section*{5 Local testing above the RM testing dimension}
In this section we look at the local characterization and local testing of $\mathrm{MRM}(q,m,d,s)$ by dimension $k$ tests, when $k$ is above the RM testing dimension $t_{q,d}=\lceil(d+1)/(q-q/p)\rceil$ of $\mathrm{RM}(q,m,d)$. By Theorem 2 dimension $k$ subspaces give a local test (and hence also a local characterization) for $\mathrm{RM}(q,m,d)$. We show they also give a local test (and hence also a local characterization) for $\mathrm{MRM}(q,m,d,s)$ for $s<q$. There is, however, some parameter loss in the reduction as we next explain. To formally state the result we need some notation. For $\boldsymbol x\in\mathbb F_q^m$ let
\begin{align*}
H_k&=\{\boldsymbol h=(\boldsymbol h_0,\boldsymbol h_1,\ldots,\boldsymbol h_k)\mid\boldsymbol h_0,\ldots,\boldsymbol h_k\in\mathbb F_q^m,\dim(\operatorname{span}\{\boldsymbol h_1,\ldots,\boldsymbol h_k\})=k\},\text{ and,}\\
H_{k,\boldsymbol x}&=\{\boldsymbol h\in H_k\mid\boldsymbol x\in\boldsymbol h_0+\operatorname{Span}\{\boldsymbol h_1,\ldots,\boldsymbol h_k\}\}.
\end{align*}
By $\boldsymbol h\sim_kH$ (resp. $H_{k,\boldsymbol x}$) we mean a choosing $\boldsymbol h$ uniformly at random from $H_k$ (resp. $H_{k,\boldsymbol x}$). We let $A_{\boldsymbol h}=\boldsymbol h_0+\operatorname{Span}\{\boldsymbol h_1,\ldots,\boldsymbol h_k\}$ be the $k$-dimensional affine space defined by $\boldsymbol h$. We also recall the distance and rejection function from the introduction:
\paragraph{Definition 50.} Let $T:\mathbb F_q^m\to\Sigma_{m,s}$ be a table. Then
\begin{itemize}
\item $\delta(T,\mathrm{MRM}(q,m,d,s))$ is the distance between $T$ and the closest evaluation table of a degree $d$ polynomial, i.e., $\delta(T,\mathrm{MRM}(q,m,d,s))=\min_{G\in\mathrm{MRM}(q,m,d,s)}\{\delta(T,G)\}$, and,
\item $\mathrm{REJ}_{k,d}^{\mathrm{MRM}}(T)$ is the probability a dimension-$k$ test demonstrates that $T$ is not a degree $d$ polynomial.
\end{itemize}
With this notation we prove:

\paragraph{Theorem 51.} Let $\mathbb F_q$ be a field of size $q$, and assume $s\leq\min\{d,q-1\}$. Suppose for $\mathrm{RM}(q,m,d)$ there exists $\alpha>0$ and $c_0\leq1$ such that for every $f$
\[\mathrm{REJ}_{k,d}^{\mathrm{RM}}(f)\geq\min\{\alpha\cdot\delta(f,\mathrm{RM}(q,m,d)),c_0\}.\]
Then, for every $s<q$, for every $T$ we have
\[\mathrm{REJ}_{k,d}^{\mathrm{MRM}}(T)\geq\min\{\alpha'\cdot\delta(T,\mathrm{MRM}(q,m,d,s)),c_0\}\tag{6}\]
for
\[\alpha'=\left(1-\frac{s-1}{q}\right)\frac1{1+q^{d/(q-1)}\frac1\alpha}.\tag{7}\]

For example, assume $q$ is prime. If $d+1\leq q-1$ the testing dimension is $1$, meaning that lines are a good local test for $\mathrm{RM}(q,m,d)$. Then, Theorem 51 says that lines are also a good test for $\mathrm{MRM}(q,m,d,s)$, alas, with a larger coefficient. Similarly, if $q<d\leq2(q-1)$ the RM testing dimension is $2$ and therefore planes are a good local test for $\mathrm{RM}(q,m,d)$. Then, Theorem 51 says that planes are also a good test for $\mathrm{MRM}(q,m,d,s)$, with a larger coefficient. In the same vein 3-dimensional planes are a good local test for $\mathrm{MRM}(q,m,d,s)$ when $d\leq3(q-1)$.
To explain the intuition behind the theorem we first consider the characterization aspect of it. Suppose $k$ is above the RM testing dimension, and that the table $T$ passes all $\mathrm{MRM}(q,m,d,s)$ $k$-dimensional tests $\boldsymbol h$. Specifically, this means that for every $\boldsymbol h\in H_k$, the table $T$ restricted to the $k$-dimensional affine space $A_{\boldsymbol h}$ is consistent with a degree $d$ polynomial $P_{\boldsymbol h}$. Now, let $T_{(0)}$ be the table $T$ where at each entry we keep only the evaluation of the function itself and remove the evaluations that are associated with higher derivatives. Then, in particular, for every $\boldsymbol h\in H_k$, the table $T_{(0)}$ restricted to the $k$-dimensional space $A_{\boldsymbol h}$ is still consistent with the degree $d$ polynomial $P_{\boldsymbol h}$. As $k$ is above the RM testing dimension, there must be a unique degree $d$ polynomial $P$ that is consistent with $P_{\boldsymbol h}$ (and the table $T_{(0)}$) for every $\boldsymbol h\in H_k$. This $P$ is the only possible candidate for a low-degree explanation of the table $T$. What we need to check, and is indeed correct, is that since $T$ passes all dimension $k$ tests, $P$ is indeed consistent with $T$.
The testing case requires more technical details but is similar in spirit. We first, again, look at $T_{(0)}$ that contains only the function evaluations, and not the higher derivatives evaluations. If $T$ passes the test with high enough probability, then so does $T_{(0)}$, and this ensures the existence of a global degree-$d$ polynomial $P$ that agrees with most values of $T_{(0)}$. Again, what remains to be shown is that $P$ agrees with most values of $T$, which we indeed prove.
In short, one can informally say that Theorem 51 shows that the MRM testing dimension is not larger than the RM testing dimension, and that above the RM testing dimension one can get local testing for MRM codes. A natural question is whether the MRM testing dimension is equal to the RM testing dimension, or not. In Section 4 we saw that when the RM testing dimension is larger than $1$, so is the MRM testing dimension, and lines do not characterize the MRM code. One might be drawn to the conjecture that the RM and MRM testing dimensions coincide. However, surprisingly, in Section 6 we show that no matter what the RM testing dimension is, when $d<sq$ the MRM testing dimension is at most two (for a precise statement see Theorem 60). For example, for $\mathrm{MRM}(q,d,m,s)$ with $2q<d<3(q-1)$, the RM testing dimension is (roughly) $3$ while the MRM testing dimension is still $2$.
\subsection*{5.1 The proof}
\paragraph{Proof.} Let $T:\mathbb F_q^m\to\Sigma_{m,s}$ be a table. Let
\[\rho=\mathrm{REJ}_{k,d}^{\mathrm{MRM}}(T)\]
be the $k$-dimensional MRM test rejection probability. If $\rho\geq c_0$ we are done. Therefore, we assume $\rho<c_0$. We first utilize what we know at the zero level.
\paragraph{Claim 53.} Let $\delta^{\mathrm{RM}}=\delta_{q,m,d}^{\mathrm{RM}}$ be the distance of the $\mathrm{RM}(q,m,d)$ code. There exists a degree $d$ polynomial $P$ such that
\[\Pr_{\boldsymbol h\in H_k}((\phi_{\boldsymbol h}\circ T)|_{A_{\boldsymbol h}}\neq\mathrm{EVAL}(P|_{A_{\boldsymbol h}}))\stackrel{\mathrm{def}}=\varepsilon_0\leq\rho\left(1+\frac1{\alpha\cdot\delta^{\mathrm{RM}}}\right).\]
We call $\boldsymbol h$ good if $(\phi_{\boldsymbol h}\circ T)|_{A_{\boldsymbol h}}=\mathrm{EVAL}(P|_{A_{\boldsymbol h}})$ and bad otherwise. In this terminology, $\varepsilon_0=\Pr_{\boldsymbol h\in H_k}[\boldsymbol h\text{ is bad}]$.
\paragraph{Proof.} Let $T_{(0)}:\mathbb F_q^m\to\mathbb F_q$ be as in Section 3, i.e., the table where we keep only the entries of the function evaluations and ignore the evaluations of higher order derivatives. For a given affine space $A$ we have that $(T|_A)_{(0)}=T_{(0)}|_A$, and so if the former agrees with a degree $d$ polynomial, so does the latter. It follows that
\[\mathrm{REJ}_{k,d}^{\mathrm{RM}}(T_{(0)})\leq\rho.\]
By the hypothesis regarding $\mathrm{RM}(q,m,d)$ and using $\rho<c_0$ there exists a unique degree $d$ polynomial $P\in\mathbb F_q[\boldsymbol X]$ such that
\[\delta_{(0)}\stackrel{\mathrm{def}}=\delta(T_{(0)},P)\leq\rho/\alpha.\]
We notice that
\begin{align*}
\mathbb E_{\boldsymbol h\sim H_k}\delta(T_{(0)}|_{A_{\boldsymbol h}},P|_{A_{\boldsymbol h}})
&=\Pr_{\boldsymbol x\sim\mathbb F_q^m}(T_{(0)}(\boldsymbol x)\neq P(\boldsymbol x))\\
&=\delta(T_{(0)},P)=\delta_{(0)},
\end{align*}
because the subspaces in $A_{\boldsymbol h}$ for $\boldsymbol h\in H_k$, cover every point an equal number of times. Therefore, by Markov's inequality,
\[\Pr_{\boldsymbol h\sim H_k}(\delta(T_{(0)}|_{A_{\boldsymbol h}},P|_{A_{\boldsymbol h}})\geq\delta^{\mathrm{RM}})\leq\frac{\delta_{(0)}}{\delta^{\mathrm{RM}}}.\tag{8}\]
Also, by assumption,
\[\Pr_{\boldsymbol h\in H_k}((\phi_{\boldsymbol h}\circ T)|_{A_{\boldsymbol h}}\notin\mathrm{MRM}(q,k,d,s))\leq\rho.\tag{9}\]
This means that except for probability $\rho$ over $\boldsymbol h$, $(\phi_{\boldsymbol h}\circ T)|_{A_{\boldsymbol h}}$ agrees with $\mathrm{EVAL}(P_{\boldsymbol h})$ for some $k$-variate degree $d$ polynomial $P_{\boldsymbol h}$. When this happens it also holds that $T_{(0)}$ agrees with $P_{\boldsymbol h}$ over $A_{\boldsymbol h}$ (because $\phi_{\boldsymbol h}$ does not change the zero level).
Thus, Equations (8) and (9) together imply that except for probability $\rho+\delta_0/\delta^{\mathrm{RM}}\leq\rho(1+1/(\alpha\delta^{\mathrm{RM}}))$ over $\boldsymbol h$, we simultaneously have that $T_{(0)}=P_{\boldsymbol h}$ over $A_{\boldsymbol h}$ and that $\delta(T_{(0)}|_{A_{\boldsymbol h}},P|_{A_{\boldsymbol h}})<\delta^{\mathrm{RM}}$. When both events happen we conclude that
\[\delta(P_{\boldsymbol h},P|_{A_{\boldsymbol h}})<\delta^{\mathrm{RM}}.\]
As $P_{\boldsymbol h}$ and $P|_{A_{\boldsymbol h}}$ are degree $d$ polynomials on $A_{\boldsymbol h}$ and are closer than $\delta^{\mathrm{RM}}$, it must be the that in fact $P_{\boldsymbol h}=P|_{A_{\boldsymbol h}}$. Thus, $(\phi_{\boldsymbol h}\circ T)|_{A_{\boldsymbol h}}$ is a valid $\mathrm{MRM}(q,k,d,s)$ table, and it is the table of the polynomial $P|_{A_{\boldsymbol h}}$, i.e., $(\phi_{\boldsymbol h}\circ T)|_{A_{\boldsymbol h}}=\mathrm{EVAL}(P|_{A_{\boldsymbol h}})$ as desired.\hfill$\triangleleft$
Our goal is to bound $\delta(\mathrm{EVAL}(P),T)$. This means that at most $\boldsymbol x\in\mathbb F_q^m$ we should have $T(\boldsymbol x)=\mathrm{EVAL}(P;\boldsymbol x)$. $T(\boldsymbol x)$ and $\mathrm{EVAL}(P;\boldsymbol x)$ contains values for all the $m$-variate directional derivatives of order up to $s$. Our handle on these values is Claim 53, that shows that most $\boldsymbol h$ are good, meaning that the $k$-variate derivatives of $P|_{A_{\boldsymbol h}}$ and $\phi_{\boldsymbol h}\circ T|_{A_{\boldsymbol h}}$ are the same. We notice that every such $k$-variate derivative is a linear combination (dependent on $\boldsymbol h$) of the $m$-variate derivatives. If $\boldsymbol x$ is such that for many $\boldsymbol h\in H_{k,\boldsymbol x}$, $\boldsymbol h$ is good, then for that $\boldsymbol x$ we get many linear equations on the $m$-variate derivatives. Our task is to prove that for many $\boldsymbol x$ there are enough good $\boldsymbol h\in H_{k,\boldsymbol x}$ to force the underlying $m$-variate derivatives of $P$ and $T$ to agree.
\paragraph{Claim 54.} We say $\boldsymbol x\in\mathbb F_q^m$ is bad if $\Pr_{\boldsymbol h\in H_{k,\boldsymbol x}}[\boldsymbol h\text{ is bad}]\geq1-(s-1)/q$ and good otherwise. Then
\[\Pr_{\boldsymbol x\in\mathbb F_q^m}[\boldsymbol x\text{ is bad}]\leq\frac q{q-(s-1)}\cdot\Pr_{\boldsymbol h\in H_k}[\boldsymbol h\text{ is bad}].\]
\paragraph{Proof.} We have
\[\mathbb E_{\boldsymbol x\in\mathbb F_q^m}\left[\Pr_{\boldsymbol h\in H_{k,\boldsymbol x}}(\boldsymbol h\text{ is bad})\right]=\Pr_{\boldsymbol h\in H_k}(\boldsymbol h\text{ is bad})=\varepsilon_0,\]
where $\varepsilon_0$ is as in Claim 53. This is because choosing a uniform $\boldsymbol h\in H_k$ is the same as first choosing a uniform $\boldsymbol x\in\mathbb F_q^m$ and then choosing a uniform $\boldsymbol h\in H_{k,\boldsymbol x}$. Therefore, for every $c>1$, by Markov's inequality,
\[\Pr_{\boldsymbol x\in\mathbb F_q^m}\left[\Pr_{\boldsymbol h\in H_{k,\boldsymbol x}}(\boldsymbol h\text{ is bad})\geq c\cdot\varepsilon_0\right]\leq1/c.\]
Choosing $c=(q-(s-1))/(q\cdot\varepsilon_0)$ gives the result.\hfill$\triangleleft$
We note that if most $k$-dimensional subspaces are good, then most lines are:
\paragraph{Claim 55.} For every $\boldsymbol x\in\mathbb F_q^m$ and $k\geq1$,
\[\Pr_{\boldsymbol u\in H_{1,\boldsymbol x}}[\boldsymbol u\text{ is bad}]\leq\Pr_{\boldsymbol h\in H_{k,\boldsymbol x}}[\boldsymbol h\text{ is bad}].\]
\paragraph{Proof.} Fix $\boldsymbol x$. For every $\boldsymbol h\in H_{k,\boldsymbol x}$, If $\boldsymbol h$ is good, then $(\phi_{\boldsymbol h}\circ T)|_{A_{\boldsymbol h}}=\mathrm{EVAL}(P|_{A_{\boldsymbol h}})$. This implies, in particular, that for every $\boldsymbol u\in H_{1,\boldsymbol x}$ such that $A_{\boldsymbol u}\subseteq A_{\boldsymbol h}$, we also have that $(\phi_{\boldsymbol u}\circ T)|_{A_{\boldsymbol u}}=\mathrm{EVAL}(P|_{A_{\boldsymbol u}})$. The result then follows because we can sample $\boldsymbol u\in H_{1,\boldsymbol x}$ by first sampling $\boldsymbol h\in A_{k,\boldsymbol x}$ and then choosing a random $\boldsymbol u\in H_{1,\boldsymbol x}$ such that $A_{\boldsymbol u}\subseteq A_{\boldsymbol h}$.\hfill$\triangleleft$
We now fix any good $\boldsymbol x$. We will prove:
\paragraph{Claim 56.} Suppose $\boldsymbol x\in\mathbb F_q^m$ is good. Then for any $m$-variate direction $\boldsymbol I$ with $w(\boldsymbol I)<s$ we have
\[T_{(\boldsymbol I)}(\boldsymbol x)=P^{(\boldsymbol I)}(\boldsymbol x).\tag{10}\]
Once we prove the claim we can conclude the proof of the theorem because:
\begin{align*}
\delta(T,\mathrm{EVAL}(P))&\leq\Pr_{\boldsymbol x\in\mathbb F_q^m}(\boldsymbol x\text{ is bad})\leq\frac q{q-(s-1)}\cdot\varepsilon_0\\
&\leq\frac q{q-(s-1)}\cdot\rho\cdot\left(1+\frac1{\alpha\cdot\delta^{\mathrm{RM}}}\right)\\
&\leq\rho\left(1-\frac q{s-1}\right)^{-1}\left(1+\frac{q^{d/(q-1)}}\alpha\right).
\end{align*}
\paragraph{Proof of Claim 56.} Fix a good $\boldsymbol x\in\mathbb F_q^m$ and let $w_0<s$. We will show Equation (10) simultaneously for all $m$-variate directions $\boldsymbol I$ with $w(\boldsymbol I)=w_0$. We know that
\[\Pr_{\boldsymbol u\in H_{1,\boldsymbol x}}(\boldsymbol u\text{ is bad})\leq\Pr_{\boldsymbol h\in H_{k,\boldsymbol x}}(\boldsymbol h\text{ is bad})<1-\frac{s-1}q,\tag{11}\]
where the first inequality is by Claim 55 and the second because $\boldsymbol x$ is good.
Suppose $\boldsymbol u=(\boldsymbol b,\boldsymbol a)$ is good (where $\boldsymbol b,\boldsymbol a\in\mathbb F_q^m$). Then, by definition, $(\phi_{\boldsymbol u}\circ T)|_{A_{\boldsymbol u}}=\mathrm{EVAL}(P|_{A_{\boldsymbol u}})$. In particular $\phi_{\boldsymbol u}(T(\boldsymbol x))=\mathrm{EVAL}(P|_{A_{\boldsymbol u}};\boldsymbol x)$, because $\boldsymbol x\in A_\ell$. Now:
\begin{itemize}
\item By Definition 48 (and plugging $k=1$)
\[(\phi_{\boldsymbol u}(T(\boldsymbol x)))_{w_0}=\sum_{\boldsymbol I,w(\boldsymbol I)=w_0}T(\boldsymbol x)_{(\boldsymbol I)}\cdot\boldsymbol a^{\boldsymbol I},\]
\item Also, by Lemma 17 (plugging $t=0$),
\[(\mathrm{EVAL}(P|_{A_{\boldsymbol u}};\boldsymbol x))_{w_0}=\sum_{\boldsymbol I,w(\boldsymbol I)=w_0}P^{\boldsymbol I}(\boldsymbol x)\cdot\boldsymbol a^{\boldsymbol I}.\]
\end{itemize}
Thus, a good $\boldsymbol u=(\boldsymbol b,\boldsymbol a)$ gives the liner equation
\[\sum_{\boldsymbol I,w(\boldsymbol I)=w_0}(P^{\boldsymbol I}(\boldsymbol x)-T(\boldsymbol x)_{(\boldsymbol I)})\cdot\boldsymbol a^{\boldsymbol I}=0,\]
where the variables are $v_{\boldsymbol I}=P^{\boldsymbol I}(\boldsymbol x)-T(\boldsymbol x)_{(\boldsymbol I)}$ for every $m$-variate direction $\boldsymbol I$ of total weight exactly $w_0$. Thus, Equation (11) implies that
\[\Pr_{\boldsymbol a\in\mathbb F_q^m\setminus\{\boldsymbol0\}}\left(\sum_{\boldsymbol I,w(\boldsymbol I)=w_0}(T(\boldsymbol x)_{(\boldsymbol I)}-P^{(\boldsymbol I)}(\boldsymbol x))\cdot\boldsymbol a^{\boldsymbol I}=0\right)>\frac{s-1}q.\]
Now, look at the polynomial $f_{\boldsymbol x}\in\mathbb F_q[X_1,\ldots,X_m]$ defined by
\[f_{\boldsymbol x}(\boldsymbol a)=\sum_{\boldsymbol I,w(\boldsymbol I)=w_0}(T(\boldsymbol x)_{(\boldsymbol I)}-P^{(\boldsymbol I)}(\boldsymbol x))\cdot\boldsymbol a^{\boldsymbol I}.\]
$f_{\boldsymbol x}$ is an $m$-variate, degree $w_0$ homogeneous polynomial, and it is $0$ with probability larger than $(s-1)/q\geq w_0/q=\deg(f_{\boldsymbol x})/q$. By the Schwartz-Zippel lemma, it must be the $0$ polynomial. Therefore, $T(\boldsymbol x)_{(\boldsymbol I)}-P^{(\boldsymbol I)}(\boldsymbol x)$ for all $\boldsymbol I$ with $w(\boldsymbol I)=w_0$ as desired.\hfill$\triangleleft$\par\hfill$\triangleleft$
\paragraph{Remark 57.} Another way to view the argument, is that each $\boldsymbol a\in\mathbb F_q^m$ is an evaluation point of a homogeneous $\mathrm{RM}(q,m,w_0)$ codeword, and therefore each good $\boldsymbol u=(\boldsymbol b,\boldsymbol a)$ gives a zero coordinate of the codeword. If the number of good $\boldsymbol u$ is too large, we get too many zeroes, and therefore the codeword must be the zero codeword, meaning that the values of the variables are zero as we wish.
\paragraph{Remark 58.} The argument we use has the information that many $k$-dimensional restrictions are good, but then chooses to reduce this knowledge to the weaker statement that for many $\boldsymbol x$, for many lines passing through $\boldsymbol x$, the linear restrictions are good. It seems that using the stronger statement might give a better code and improve the parameters, but we have not succeeded yet in analyzing this.

\begin{thebibliography}{20}
\bibitem[1]{ref1} Noga Alon. Combinatorial nullstellensatz. Combinatorics, Probability and Computing, 8(1- 2):7–29, 1999.
\bibitem[2]{ref2} Noga Alon, Tali Kaufman, Michael Krivelevich, Simon Litsyn, and Dana Ron. Testing reed-muller codes. IEEE Transactions on Information Theory, 51(11):4032–4039, 2005.
\bibitem[3]{ref3} Simeon Ball and Oriol Serra. Punctured combinatorial nullstellensätze. Combinatorica, 29(5):511–522, 2009.
\bibitem[4]{ref4} Arnab Bhattacharyya, Swastik Kopparty, Grant Schoenebeck, Madhu Sudan, and David Zuckerman. Optimal testing of reed-muller codes. In 2010 IEEE 51st Annual Symposium on Foundations of Computer Science, pages 488–497. IEEE, 2010.
\bibitem[5]{ref5} David Cox, John Little, and Donal OShea. Ideals, varieties, and algorithms: an introduction to computational algebraic geometry and commutative algebra. Springer Science \& Business Media, 2013.
\bibitem[6]{ref6} Zeev Dvir, Swastik Kopparty, Shubhangi Saraf, and Madhu Sudan. Extensions to the method of multiplicities, with applications to kakeya sets and mergers. SIAM Journal on Computing, 42(6):2305–2328, 2013.
\bibitem[7]{ref7} Katalin Friedl and Madhu Sudan. Some improvements to total degree tests. In Proceedings Third Israel Symposium on the Theory of Computing and Systems, pages 190–198. IEEE, 1995.
\bibitem[8]{ref8} Peter Gemmell, Richard J. Lipton, Ronitt Rubinfeld, Madhu Sudan, and Avi Wigderson. Self-testing/correcting for polynomials and for approximate functions. In Cris Koutsougeras and Jeffrey Scott Vitter, editors, Proceedings of the 23rd Annual ACM Symposium on Theory of Computing, May 5-8, 1991, New Orleans, Louisiana, USA, pages 32–42. ACM, 1991. doi:10.1145/103418.103429.
\bibitem[9]{ref9} Venkatesan Guruswami and Carol Wang. Linear-algebraic list decoding for variants of reed– solomon codes. IEEE Transactions on Information Theory, 59(6):3257–3268, 2013.
\bibitem[10]{ref10} Elad Haramaty, Amir Shpilka, and Madhu Sudan. Optimal testing of multivariate polynomials over small prime fields. SIAM Journal on Computing, 42(2):536–562, 2013.
\bibitem[11]{ref11} Charanjit S Jutla, Anindya C Patthak, Atri Rudra, and David Zuckerman. Testing low-degree polynomials over prime fields. In 45th Annual IEEE Symposium on Foundations of Computer Science, pages 423–432. IEEE, 2004.
\bibitem[12]{ref12} Tali Kaufman and Dana Ron. Testing polynomials over general fields. SIAM Journal on Computing, 36(3):779–802, 2006.
\bibitem[13]{ref13} Swastik Kopparty. Some remarks on multiplicity codes. Discrete Geometry and Algebraic Combinatorics, 625:155–176, 2013.
\bibitem[14]{ref14} Swastik Kopparty, Or Meir, Noga Ron-Zewi, and Shubhangi Saraf. High-rate locally correctable and locally testable codes with sub-polynomial query complexity. Journal of the ACM (JACM), 64(2):1–42, 2017.
\bibitem[15]{ref15} Swastik Kopparty, Shubhangi Saraf, and Sergey Yekhanin. High-rate codes with sublinear-time decoding. Journal of the ACM (JACM), 61(5):1–20, 2014.
\bibitem[16]{ref16} Rasmus Refslund Nielsen. List decoding of linear block codes. PhD thesis, DTU, 2001.
\bibitem[17]{ref17} M Yu Rosenbloom and Michael Anatol’evich Tsfasman. Codes for the m-metric. Problemy Peredachi Informatsii, 33(1):55–63, 1997.
\bibitem[18]{ref18} Ronitt Rubinfeld and Madhu Sudan. Self-testing polynomial functions efficiently and over rational domains. In Greg N. Frederickson, editor, Proceedings of the Third Annual ACM/SIGACT-SIAM Symposium on Discrete Algorithms, 27-29 January 1992, Orlando, Florida, USA, pages 23–32. ACM/SIAM, 1992. URL: http://dl.acm.org/citation.cfm?id=139404.139410.
\bibitem[19]{ref19} Ronitt Rubinfeld and Madhu Sudan. Robust characterizations of polynomials with applications to program testing. SIAM Journal on Computing, 25(2):252–271, 1996.
\bibitem[20]{ref20} Madhu Sudan, Luca Trevisan, and Salil Vadhan. Pseudorandom generators without the xor lemma. Journal of Computer and System Sciences, 62(2):236–266, 2001.
\end{thebibliography}
\end{document}
